\documentclass[11pt, a4paper]{article}

% ── PACKAGES ──────────────────────────────────────────────────
\usepackage[utf8]{inputenc}
\usepackage[T1]{fontenc}
\usepackage{geometry}
\geometry{margin=2.2cm}
\usepackage{hyperref}
\hypersetup{
    colorlinks=true,
    linkcolor=blue,
    urlcolor=blue,
    citecolor=blue,
    pdftitle={FOXA1/EZH2 Dual IHC Protocol Specification:
              Antibody Selection, Staining Conditions,
              Scoring Method, and Calibration Study Design},
    pdfauthor={Eric Robert Lawson / OrganismCore}
}
\usepackage{amsmath}
\usepackage{booktabs}
\usepackage{array}
\usepackage{longtable}
\usepackage{parskip}
\usepackage{microtype}
\usepackage{authblk}
\usepackage{titlesec}
\usepackage{fancyhdr}
\usepackage{xcolor}
\usepackage{float}
\usepackage{enumitem}
\usepackage{multirow}
\usepackage{mdframed}

% ── COLOURS ───────────────────────────────────────────────────
\definecolor{repobox}{RGB}{240,247,255}
\definecolor{protocolbox}{RGB}{245,255,245}
\definecolor{warningbox}{RGB}{255,248,210}
\definecolor{gapbox}{RGB}{255,235,235}
\definecolor{scorebox}{RGB}{250,245,255}
\definecolor{decisionbox}{RGB}{235,255,235}
\definecolor{headerblue}{RGB}{30,80,160}

% ── HEADER / FOOTER ───────────────────────────────────────────
\pagestyle{fancy}
\fancyhf{}
\lhead{\footnotesize FOXA1/EZH2 IHC Protocol Specification
       --- OrganismCore}
\rhead{\footnotesize Version 1.0 \textbar{} 2026-03-06}
\rfoot{\small Page \thepage}
\renewcommand{\headrulewidth}{0.4pt}
\setlength{\headheight}{24pt}
\addtolength{\topmargin}{-10pt}

% ── SECTION FORMATTING ────────────────────────────────────────
\titleformat{\section}
  {\large\bfseries\color{headerblue}}
  {\thesection}{1em}{}[\titlerule]
\titleformat{\subsection}
  {\normalsize\bfseries}{\thesubsection}{1em}{}
\titleformat{\subsubsection}
  {\small\bfseries\itshape}{\thesubsubsection}{1em}{}

% ══════════════════════════════════════════════════════════════
\title{
    \vspace{-1em}
    \textbf{FOXA1/EZH2 Dual IHC Protocol Specification}\\[0.5em]
    \large Antibody Selection, Staining Conditions,
    H-Score Calculation, Tissue Controls,
    and Calibration Study Design\\[0.4em]
    \normalsize Supporting document for
    CS-LIT-1 (\href{https://doi.org/10.5281/zenodo.18883922}
    {DOI: 10.5281/zenodo.18883922}) and
    CS-LIT-22 (\href{https://doi.org/10.5281/zenodo.18892788}
    {DOI: 10.5281/zenodo.18892788})\\[0.3em]
    \small Version 1.0 --- Pre-calibration specification.
    Cut-points to be determined in the concordance study
    described in Section~\ref{sec:calibration}.
}

\author[1]{Eric Robert Lawson}
\affil[1]{OrganismCore \textbar{}
    \href{mailto:OrganismCore@proton.me}
    {OrganismCore@proton.me} \textbar{}
    ORCID: 0009-0002-0414-6544}

\date{March 6, 2026}

% ══════════════════════════════════════════════════════════════
\begin{document}
\maketitle
\thispagestyle{fancy}

% ── VERSION NOTE ──────────────────────────────────────────────
\begin{mdframed}[backgroundcolor=warningbox,
                 linecolor=orange!70!black,
                 linewidth=1pt]
\textbf{VERSION NOTE:} This is a pre-calibration
protocol specification. Antibody recommendations
and staining conditions are based on published
breast cancer IHC literature and manufacturer
datasheets. The H-score cut-points for subtype
classification (Section~\ref{sec:cutpoints}) are
provisional estimates derived from scRNA-seq
ratio values and require calibration against
PAM50-classified FFPE tissue before clinical
deployment. See Section~\ref{sec:calibration}
for the calibration study design.
\end{mdframed}

\vspace{0.5em}

\tableofcontents

\newpage

% ══���═══════════════════════════════════════════════════════════
\section{Purpose and Scope}
% ══════════════════════════════════════════════════════════════

This document provides the complete technical
specification for running the FOXA1/EZH2 dual IHC
ratio assay on formalin-fixed, paraffin-embedded
(FFPE) breast cancer tissue. It is intended for
use by:

\begin{itemize}
    \item Pathologists and histotechnologists
    implementing the assay in a clinical or
    research setting
    \item Researchers designing the IHC
    concordance study described in CS-LIT-22
    \item Institutions evaluating the assay
    for inclusion in a breast cancer
    diagnostic panel
\end{itemize}

The document covers:
\begin{enumerate}
    \item Antibody clone selection with catalogue
    references and manufacturer datasheet links
    \item Staining protocol for manual and
    automated platforms
    \item Tissue section preparation
    \item Scoring method and H-score calculation
    \item Provisional cut-points and the
    calibration study needed to confirm them
    \item Positive and negative tissue controls
    \item Troubleshooting guidance
\end{enumerate}

% ══════════════════════════════════════════════════════════════
\section{Antibody Specifications}
% ══════════════════════════════════════════════════════════════

\subsection{FOXA1 Antibody}

\begin{center}
\begin{tabular}{p{4.5cm}p{9cm}}
\toprule
\textbf{Field} & \textbf{Specification} \\
\midrule
Target protein & FOXA1 (Forkhead Box A1;
    Hepatocyte Nuclear Factor 3-alpha) \\[0.3em]
\textbf{Recommended clone} &
    \textbf{D4E2 (XP Rabbit Monoclonal)} \\[0.3em]
Manufacturer & Cell Signaling Technology \\[0.3em]
Catalogue number & CST \#53528 \\[0.3em]
Host species & Rabbit \\[0.3em]
Clonality & Monoclonal \\[0.3em]
Immunogen & Synthetic peptide corresponding
    to human FOXA1 \\[0.3em]
Molecular weight & ${\sim}49$~kDa \\[0.3em]
Cellular localisation & Nuclear (exclusively) \\[0.3em]
Validated for FFPE IHC & Yes (manufacturer
    validated) \\[0.3em]
Manufacturer IHC datasheet &
    \href{https://www.cellsignal.com/products/primary-antibodies/foxa1-d4e2-xp-rabbit-mab/53528}
    {cellsignal.com \#53528} \\[0.3em]
\textbf{Alternative clone} &
    SP87 (Rabbit Monoclonal) ---
    Abcam ab170933;
    equivalent performance in
    published FOXA1 IHC studies \\[0.3em]
Why this antibody & D4E2 is validated in
    breast cancer IHC by the manufacturer
    and by multiple published studies.
    Nuclear-only staining pattern simplifies
    scoring. \\
\bottomrule
\end{tabular}
\end{center}

\subsection{EZH2 Antibody}

\begin{center}
\begin{tabular}{p{4.5cm}p{9cm}}
\toprule
\textbf{Field} & \textbf{Specification} \\
\midrule
Target protein & EZH2 (Enhancer of Zeste
    Homolog 2; KMT6) \\[0.3em]
\textbf{Recommended clone} &
    \textbf{D2H1 (XP Rabbit Monoclonal)} \\[0.3em]
Manufacturer & Cell Signaling Technology \\[0.3em]
Catalogue number & CST \#5246 \\[0.3em]
Host species & Rabbit \\[0.3em]
Clonality & Monoclonal \\[0.3em]
Immunogen & Synthetic peptide corresponding
    to human EZH2 (around residue 700) \\[0.3em]
Molecular weight & ${\sim}85$~kDa \\[0.3em]
Cellular localisation & Nuclear (exclusively) \\[0.3em]
Validated for FFPE IHC & Yes (manufacturer
    validated) \\[0.3em]
Manufacturer IHC datasheet &
    \href{https://www.cellsignal.com/products/primary-antibodies/ezh2-d2h1-xp-rabbit-mab/5246}
    {cellsignal.com \#5246} \\[0.3em]
\textbf{Alternative clone} &
    AC22 (Mouse Monoclonal) ---
    Leica Biosystems NCL-L-EZH2;
    widely used in published breast
    cancer EZH2 IHC studies;
    requires anti-mouse secondary \\[0.3em]
Why this antibody & D2H1 is the most
    extensively published EZH2 IHC
    antibody in breast cancer.
    Nuclear-only staining pattern matches
    FOXA1, enabling consistent
    H-score comparison on adjacent
    or sequential sections. \\
\bottomrule
\end{tabular}
\end{center}

\begin{mdframed}[backgroundcolor=protocolbox,
                 linecolor=green!50!black,
                 linewidth=1pt]
\textbf{NOTE ON ANTIBODY MATCHING:}
Both D4E2 (FOXA1) and D2H1 (EZH2) are rabbit
monoclonal antibodies from the same manufacturer
(Cell Signaling Technology). This means:
\begin{itemize}[noitemsep]
    \item The same secondary antibody (anti-rabbit
    HRP polymer) can be used for both stains
    when run sequentially
    \item Protocol conditions (blocking, incubation,
    detection) are directly comparable
    \item Both produce exclusively nuclear staining,
    simplifying H-score calculation to nuclear
    compartment only
\end{itemize}
\end{mdframed}

% ══════════════════════════════════════════════════════════════
\section{Tissue Section Preparation}
% ══════════════════════════════════════════════════════════════

\subsection{Section thickness and preparation}

\begin{itemize}
    \item \textbf{Section thickness:} 3--4~$\mu$m
    (standard diagnostic IHC thickness)
    \item \textbf{Sections required per case:}
    Two serial or near-serial sections ---
    one for FOXA1, one for EZH2. Sections should
    be taken within 10 levels of each other to
    ensure the same tumour area is represented
    in both stains.
    \item \textbf{Mounting:} Charged or
    poly-L-lysine coated slides (APES-coated
    slides acceptable) to prevent section
    detachment during antigen retrieval
    \item \textbf{Drying:} 60$^\circ$C oven,
    minimum 60 minutes, after mounting
    \item \textbf{Tumour content:} Select
    section with $\geq 20\%$ invasive tumour
    cells. A pathologist should mark the tumour
    area on the H\&E before IHC scoring.
\end{itemize}

\subsection{FFPE block age and RNA quality}

IHC does not require RNA. FFPE blocks of any
age are suitable provided:
\begin{itemize}[noitemsep]
    \item Standard formalin fixation was used
    (10\% neutral buffered formalin, 6--72 hours)
    \item The block has not undergone repeated
    freeze-thaw cycling
    \item Section quality is adequate (no tears,
    folds, or edge artefacts in tumour area)
\end{itemize}

This is a significant advantage over PAM50/Prosigna,
which requires adequate RNA integrity and is
unreliable on poorly fixed or archival tissue.

% ══════════════════════════════════════════════════════════════
\section{Staining Protocol}
% ══════════════════════════════════════════════════════════════

\subsection{Manual protocol (all platforms)}

\begin{center}
\begin{tabular}{p{0.8cm}p{5cm}p{7.5cm}}
\toprule
\textbf{Step} & \textbf{Procedure} &
\textbf{Conditions / Notes} \\
\midrule
1 & Deparaffinise &
    Xylene, $3 \times 5$ min; or equivalent
    deparaffinisation solution \\[0.3em]
2 & Rehydrate &
    Graded ethanol series: 100\%, 95\%, 70\%,
    distilled water, $2$ min each \\[0.3em]
3 & Endogenous peroxidase block &
    3\% $\text{H}_2\text{O}_2$ in methanol
    or water, 10 min, room temperature (RT).
    Wash $3\times$ PBS. \\[0.3em]
4 & \textbf{Antigen retrieval ---
    FOXA1 (D4E2)} &
    \textbf{HIER, Citrate buffer pH~6.0
    (e.g. Dako S1699).}
    Pressure cooker or steamer, 20 min at
    $95$--$100^\circ$C.
    Cool 20--30 min at RT.
    Wash $2\times$ PBS. \\[0.3em]
4 & \textbf{Antigen retrieval ---
    EZH2 (D2H1)} &
    \textbf{HIER, EDTA buffer pH~9.0
    (e.g. Dako S2367 or Leica AR9640).}
    Pressure cooker or steamer, 20 min at
    $98^\circ$C.
    Cool 20--30 min at RT.
    Wash $2\times$ PBS.
    \textit{Note: run FOXA1 and EZH2
    sections separately for antigen
    retrieval due to different buffer
    requirements.} \\[0.3em]
5 & Serum block &
    5\% normal goat serum (or equivalent
    protein block) in PBS, 30 min, RT.
    Do not wash. Tip off excess. \\[0.3em]
6 & \textbf{Primary antibody ---
    FOXA1} &
    \textbf{D4E2 (CST \#53528) at 1:100
    in antibody diluent (e.g. Dako S0809).}
    60 min at RT in humidified chamber.
    Wash $3\times$ PBS, 5 min each. \\[0.3em]
6 & \textbf{Primary antibody ---
    EZH2} &
    \textbf{D2H1 (CST \#5246) at 1:100
    in antibody diluent.}
    60 min at RT in humidified chamber.
    Wash $3\times$ PBS, 5 min each. \\[0.3em]
7 & Secondary antibody &
    Anti-rabbit HRP-polymer
    (e.g. Dako EnVision$+$ System HRP,
    Rabbit, K4003; or equivalent).
    30 min at RT.
    Wash $3\times$ PBS. \\[0.3em]
8 & Chromogen development &
    DAB (3,3'-diaminobenzidine)
    substrate-chromogen solution.
    3--8 min at RT.
    Monitor microscopically.
    Stop with distilled water when
    brown nuclear staining is visible
    and background is minimal. \\[0.3em]
9 & Counterstain &
    Mayer's haematoxylin, 2 min.
    Blueing in running tap water, 5 min. \\[0.3em]
10 & Dehydrate and mount &
    Graded ethanol, xylene, permanent
    mounting medium (e.g. DPX). \\
\bottomrule
\end{tabular}
\end{center}

\subsection{Automated platform adaptations}

\begin{center}
\begin{tabular}{p{3.5cm}p{5cm}p{5cm}}
\toprule
\textbf{Platform} & \textbf{FOXA1 (D4E2)} &
\textbf{EZH2 (D2H1)} \\
\midrule
Ventana BenchMark
(Roche) &
CC1 (Tris-EDTA pH~8.5) standard,
36 min at $100^\circ$C.
D4E2 diluted 1:100 in Ventana
antibody diluent.
OptiView DAB detection. &
CC1 standard, 36 min $100^\circ$C.
D2H1 diluted 1:100.
OptiView DAB detection. \\[0.4em]
Dako AutoStainer
(Agilent) &
EnVision Flex target retrieval
low pH (citrate equivalent),
20 min $98^\circ$C.
D4E2 1:100 in Flex antibody
diluent. &
EnVision Flex target retrieval
high pH (EDTA equivalent),
20 min $98^\circ$C.
D2H1 1:100. \\[0.4em]
Leica Bond RX &
Epitope Retrieval Solution 1
(citrate pH 6), ER1, 20 min.
D4E2 1:100.
Bond Polymer Refine detection. &
Epitope Retrieval Solution 2
(EDTA pH 9), ER2, 20 min.
D2H1 1:100.
Bond Polymer Refine detection. \\
\bottomrule
\end{tabular}
\end{center}

% ══════════════════════════════════════════════════════════════
\section{Tissue Controls}
% ══════════════════════════════════════════════════════════════

\begin{center}
\begin{tabular}{p{3cm}p{4.5cm}p{4.5cm}}
\toprule
\textbf{Control type} & \textbf{FOXA1} &
\textbf{EZH2} \\
\midrule
\textbf{Positive tissue control} &
Luminal A breast carcinoma
(known ER$+$, LumA by PAM50 or
clinical behaviour).
Should show strong diffuse
nuclear staining. &
TNBC or basal-like breast carcinoma
(known high EZH2 by prior IHC
or literature).
Should show strong diffuse
nuclear staining. \\[0.4em]
\textbf{Negative tissue control} &
Claudin-low or TNBC breast
carcinoma (FOXA1-absent).
Should show absent or
near-absent nuclear staining. &
Normal breast stroma.
Should show weak or absent
nuclear staining. \\[0.4em]
\textbf{Internal positive control} &
Normal ductal epithelium
on same section.
Luminal cells should show
moderate FOXA1 staining. &
Mitotically active cells
(inflammatory cells, proliferating
stroma) often show EZH2 staining;
use as process control only. \\[0.4em]
\textbf{Negative reagent control} &
Omit primary antibody; use
isotype-matched IgG at same
concentration. &
Omit primary antibody; use
isotype-matched IgG at same
concentration. \\
\bottomrule
\end{tabular}
\end{center}

% ══════════════════════════════════════════════════════════════
\section{Scoring Method}
\label{sec:scoring}
% ══════════════════════════════════════════════════════════════

\subsection{What to score}

\begin{mdframed}[backgroundcolor=scorebox,
                 linecolor=purple!50!black,
                 linewidth=1pt]
\textbf{SCORE INVASIVE CANCER CELLS ONLY.}
Exclude:
\begin{itemize}[noitemsep]
    \item Normal ductal and lobular epithelium
    \item Stromal fibroblasts and endothelial cells
    \item Tumour-infiltrating lymphocytes (TILs)
    \item Myoepithelial cells
    \item Necrotic areas
\end{itemize}
A pathologist should mark the invasive tumour area
on the H\&E before IHC scoring. Score the marked
area only. This is critical: the ratio derives its
subtype-discriminating power from cancer cell
expression, and stromal/immune cell EZH2 expression
would inflate the denominator and distort the ratio.
\end{mdframed}

\subsection{H-score calculation}

The H-score is a semi-quantitative scoring system
that accounts for both staining intensity and
proportion of positive cells. It is computed as:

\begin{equation}
\text{H-score} = \sum_{i=1}^{3} i \cdot p_i
\end{equation}

where:
\begin{itemize}
    \item $p_1$ = percentage of invasive cancer
    cells with \textbf{weak} (1$+$) nuclear staining
    \item $p_2$ = percentage of invasive cancer
    cells with \textbf{moderate} (2$+$) nuclear
    staining
    \item $p_3$ = percentage of invasive cancer
    cells with \textbf{strong} (3$+$) nuclear
    staining
    \item Percentages are estimated to the nearest
    5\% and must sum to $\leq 100\%$ (unstained
    cells are not counted in the sum)
\end{itemize}

\textbf{H-score range: 0--300.}

\begin{center}
\begin{tabular}{p{3cm}p{5cm}p{5cm}}
\toprule
\textbf{Staining
intensity} &
\textbf{Visual description} &
\textbf{Practical example} \\
\midrule
0 (negative) &
No discernible nuclear staining
above background &
FOXA1 in claudin-low; EZH2 in
normal stroma \\[0.3em]
1$+$ (weak) &
Light brown nuclear tinge;
identifiable but faint;
clearly above background &
FOXA1 in TNBC (residual low
expression) \\[0.3em]
2$+$ (moderate) &
Clearly brown nuclear staining;
unambiguous; not intense &
FOXA1 in LumB; EZH2 in LumA \\[0.3em]
3$+$ (strong) &
Intense brown nuclear staining;
immediately evident at low power &
FOXA1 in LumA; EZH2 in TNBC \\
\bottomrule
\end{tabular}
\end{center}

\subsection{Example H-score calculation}

A TNBC tumour section shows:
\begin{itemize}
    \item FOXA1: 5\% cells at 1$+$, 2\% at 2$+$,
    0\% at 3$+$
    \item $\text{FOXA1 H-score} = (1 \times 5) +
    (2 \times 2) + (3 \times 0) = 9$
\end{itemize}
\begin{itemize}
    \item EZH2: 20\% cells at 1$+$, 45\% at 2$+$,
    30\% at 3$+$
    \item $\text{EZH2 H-score} = (1 \times 20) +
    (2 \times 45) + (3 \times 30) = 20 + 90 + 90 = 200$
\end{itemize}

\textbf{FOXA1/EZH2 ratio} $= 9 / 200 = 0.045$

This falls below 0.2, consistent with the
TNBC/CL range in the predicted framework.

\subsection{Minimum tumour cell count}

Score a minimum of \textbf{200 invasive cancer
cells} for each stain. For small biopsies where
$<200$ cells are available, score all visible
invasive cells and note the count.

\subsection{Inter-observer assessment}

For the calibration study, each slide should be
independently scored by \textbf{two pathologists}
with no prior knowledge of the other's score.
Inter-observer reliability should be calculated
as:
\begin{itemize}
    \item Intraclass correlation coefficient (ICC)
    for continuous H-score values
    \item Cohen's $\kappa$ for categorical subtype
    assignment
    \item Target: ICC $> 0.80$, $\kappa > 0.70$
    before clinical deployment
\end{itemize}

% ══════════════════════════════════════════════════════════════
\section{Provisional Cut-Points}
\label{sec:cutpoints}
% ══════════════════════════════════════════════════════════════

\begin{mdframed}[backgroundcolor=warningbox,
                 linecolor=orange!70!black,
                 linewidth=1pt]
\textbf{IMPORTANT:} The following cut-points
are provisional estimates. They are derived by
mapping the scRNA-seq ratio values (FOXA1/EZH2
mean expression per subtype) to approximate
H-score ratio equivalents based on known
expression patterns in published IHC literature.
\textbf{They must be calibrated against
PAM50-classified FFPE tissue before clinical use.}
See Section~\ref{sec:calibration} for the
calibration study design.
\end{mdframed}

\begin{center}
\colorbox{decisionbox}{
\begin{minipage}{0.92\textwidth}
\vspace{0.8em}
\textbf{FOXA1/EZH2 IHC H-SCORE RATIO ---
PROVISIONAL DECISION TABLE}\\[0.6em]
\begin{tabular}{p{3cm}p{2.5cm}p{6.5cm}}
\toprule
\textbf{Subtype} &
\textbf{Provisional ratio range} &
\textbf{Treatment logic
(from framework)} \\
\midrule
Luminal A &
$> 1.5$ &
ET engages directly.
CDK4/6 inhibitor (palbociclib,
ribociclib, abemaciclib) + ET
is the primary entry. \\[0.4em]
Luminal B &
$0.8$--$1.5$ &
Chromatin lock on ER output.
HDACi (entinostat) + ET.
CDK4/6i also active. \\[0.4em]
HER2-enriched &
$0.3$--$0.8$ &
Anti-HER2 first (trastuzumab
$\pm$ pertuzumab).
ET thereafter if HR$+$. \\[0.4em]
TNBC &
$0.02$--$0.15$ &
Epigenetic lock.
Tazemetostat $\rightarrow$
fulvestrant (see CS-LIT-16). \\[0.4em]
Claudin-low &
$< 0.02$ &
No epigenetic lock to dissolve.
Anti-TIGIT $\rightarrow$
anti-PD-1 sequence required. \\[0.4em]
ILC
(exception) &
$> 2.0$ and
CDH1-absent &
FOXA1 hyperactivated above LumA.
Fulvestrant $>$ AI.
Full ER degradation required. \\
\bottomrule
\end{tabular}
\vspace{0.5em}
\end{minipage}}
\end{center}

\vspace{0.5em}

\textbf{Note on ILC identification:} ILC sits
above LumA on the ratio (FOXA1 hyperactivated,
EZH2 low). It is identified by three features
together: (1) ratio $>2.0$; (2) CDH1-absent
by IHC or clinical/morphological confirmation
of lobular histology; (3) characteristic
single-file infiltration pattern on H\&E.
CDH1 IHC can be added as a third stain
if lobular vs ductal distinction is required.

% ══════════════════════════════════════════════════════════════
\section{Calibration Study Design}
\label{sec:calibration}
% ═══════════════════��══════════════════════════════════════════

This section describes the single study needed
to translate the provisional cut-points into
clinically validated IHC thresholds.

\subsection{Study type}

\textbf{Retrospective IHC concordance study.}
No new patient recruitment. No treatment
changes. Observational only.

\subsection{Study population}

Archived FFPE breast cancer tissue from patients
with:
\begin{itemize}[noitemsep]
    \item Available PAM50/Prosigna result
    (from a prior clinical or research assay)
    OR available PAM50 result from a freshly
    run parallel assay on the same tissue
    \item Confirmed invasive breast carcinoma
    (any subtype)
    \item Adequate tumour content in FFPE
    block ($\geq 20\%$ invasive tumour cells)
\end{itemize}

\textbf{Target sample size:}
$n = 300$--$400$ cases minimum, with
representation from all five subtypes:

\begin{center}
\begin{tabular}{lcc}
\toprule
\textbf{PAM50 subtype} &
\textbf{Target n} &
\textbf{Minimum acceptable} \\
\midrule
Luminal A & 80--100 & 50 \\
Luminal B & 80--100 & 50 \\
HER2-enriched & 50--60 & 30 \\
Basal/TNBC & 60--80 & 40 \\
Claudin-low & 30--40 & 20 \\
\textbf{Total} & \textbf{300--380} & \textbf{190} \\
\bottomrule
\end{tabular}
\end{center}

\textbf{ILC cases:} Include a minimum of
30 lobular carcinoma cases (CDH1-absent, any
PAM50 subtype) as a separate analysis arm.

\subsection{Primary endpoint}

Agreement between FOXA1/EZH2 IHC ratio
classification and PAM50 subtype classification,
measured by:
\begin{itemize}
    \item \textbf{Cohen's $\kappa$:} Target $\geq 0.70$
    \item \textbf{Overall concordance:}
    Target $\geq 75\%$ (ideally $\geq 85\%$)
    \item \textbf{AUC for binary classification tasks:}
    LumA vs Basal, LumA vs LumB
    (pre-specified; target AUC $\geq 0.75$)
\end{itemize}

\subsection{Secondary endpoints}

\begin{itemize}
    \item Inter-observer reliability between
    two independent pathologists
    (ICC for H-scores; $\kappa$ for categorical
    assignment; target ICC $> 0.80$)
    \item Calibrated cut-point values in IHC
    H-score units (to replace the provisional
    cut-points in Section~\ref{sec:cutpoints})
    \item Correlation between IHC ratio and
    clinical outcome where outcome data are
    available (exploratory; not a primary endpoint)
\end{itemize}

\subsection{What this study does not require}

\begin{itemize}[noitemsep]
    \item No new patient recruitment
    \item No treatment changes to any patient
    \item No regulatory approval beyond
    standard institutional ethics
    for tissue research (most jurisdictions
    classify this as minimal-risk retrospective
    tissue research)
    \item No new laboratory equipment
    \item No bioinformatics infrastructure
\end{itemize}

\subsection{Estimated resource requirements}

\begin{center}
\begin{tabular}{p{5cm}p{7.5cm}}
\toprule
\textbf{Resource} & \textbf{Estimate} \\
\midrule
Antibody cost &
FOXA1 D4E2 (CST \#53528): ${\sim}$\$400
for 100 tests per vial. EZH2 D2H1
(CST \#5246): ${\sim}$\$400 per vial.
For 400 cases: approximately \$1,500--\$3,000
total in antibody costs. \\[0.3em]
IHC consumables &
Slides, coverslips, DAB, haematoxylin,
buffers: approximately \$2,000--\$4,000
for 400 cases. \\[0.3em]
Pathologist time &
Scoring 400 $\times$ 2 slides $\times$ 2
pathologists: approximately 40--60 hours
total. \\[0.3em]
Statistician time &
Concordance analysis, ROC, cut-point
determination: approximately 10--20 hours. \\[0.3em]
\textbf{Total estimated cost} &
\textbf{\$5,000--\$15,000 depending on
institution and existing PAM50 data
availability.} \\
\bottomrule
\end{tabular}
\end{center}

% ══════════════════════════════════════════════════════════════
\section{Troubleshooting}
% ══════════════════════════════════════════════════════════════

\begin{center}
\begin{tabular}{p{4cm}p{4cm}p{5cm}}
\toprule
\textbf{Problem} &
\textbf{Likely cause} &
\textbf{Solution} \\
\midrule
No or very weak staining
(both antibodies) &
Antigen retrieval failure;
antibody inactivation;
inadequate section adhesion &
Increase retrieval time.
Check antibody storage ($-20^\circ$C).
Use freshly cut sections ($<6$ weeks). \\[0.4em]
High background on EZH2 &
Insufficient peroxidase block;
non-specific binding &
Extend $\text{H}_2\text{O}_2$
block to 15 min.
Increase serum block concentration. \\[0.4em]
FOXA1 cytoplasmic staining &
Non-specific binding;
cross-reactivity &
Only count nuclear staining.
Cytoplasmic FOXA1 is not scored
for this ratio. \\[0.4em]
CL ratio appears higher than
TNBC (inconsistent with framework) &
Stromal cells being scored;
EZH2 in TILs included in count &
Re-score cancer cells only.
Exclude all non-cancer nuclei.
This is the most common scoring error. \\[0.4em]
Poor section quality (tears,
folds, loss of antigen) &
Block quality; section age;
inadequate drying &
Re-cut from block.
Use sections cut within 2 weeks. \\
\bottomrule
\end{tabular}
\end{center}

% ══════════════════════════════════════════════════════════════
\section{Document Status and Contact}
% ══════════════════════════════════════════════════════════════

\begin{center}
\begin{tabular}{ll}
\toprule
\textbf{Field} & \textbf{Value} \\
\midrule
Document type &
IHC Protocol Specification v1.0 \\
Status &
Pre-calibration. Cut-points provisional. \\
Author &
Eric Robert Lawson, OrganismCore \\
Contact &
OrganismCore@proton.me \\
ORCID &
0009-0002-0414-6544 \\
CS-LIT-1 DOI &
\href{https://doi.org/10.5281/zenodo.18883922}
{10.5281/zenodo.18883922} \\
CS-LIT-22 DOI &
\href{https://doi.org/10.5281/zenodo.18892788}
{10.5281/zenodo.18892788} \\
Repository &
\href{https://github.com/Eric-Robert-Lawson/attractor-oncology}
{Eric-Robert-Lawson/attractor-oncology} \\
Version history &
v1.0 (2026-03-06): Initial release \\
Next version trigger &
Publication of calibration study results;
update cut-points to validated IHC values \\
\bottomrule
\end{tabular}
\end{center}

\end{document}
